% TMLR version: condensed Sections 2.5-2.7. Full arXiv-v1 text in Appendix I (Extended methods).
\subsection{Ridge readout}
\label{sec:ridge}

Encoding models used dual ridge regression \citep{hoerl1970} on the Conv1 maps, pooled to a
$g \times g$ grid by adaptive average pooling, $g \in \{1, 2, 4, 8\}$, giving $32g^2$ features
($g = 1$ is GAP), standardised with statistics of the training rows of each split only.
Cross-validation is nested and concept-grouped (8 outer and 5 inner folds, seed 20260818); per
voxel and outer fold, $\alpha$ was chosen from 29 log-spaced values between $10^{-6}$ and $10^{8}$
by the minimum inner-fold mean squared error, not by Pearson $r$, the reported score
(Appendix~\ref{app:ridge}). Out-of-fold predictions were pooled over the outer folds and scored by
one Pearson $r$ per voxel over all 720 stimuli; the score is the median $r$ over reliable V1 voxels,
averaged over subjects, with contrasts formed within seed. Under this pooled score, features
without signal yield negative rather than zero scores.

The pre-specified analysis used the 720 single-trial betas at NATIVE 224\,px, all checkpoints and
subjects, and all four $g$; an exploratory fit at $g = 1$ covered all NATIVE resolutions and
UPSAMPLED 224. Controls with the identical readout were low-level banks (pixels; a
luminance-and-contrast bank of $2g^2$ features; a Gabor bank), Backprop seed-0 features with
permuted stimulus assignment as negative test, and an image-wise permutation of the $g \times g$
cells (seed-0 untrained and Backprop, $g \in \{2, 4, 8\}$; details in Appendix~\ref{app:ridge}).

\subsection{Mean-profile components and decomposition}
\label{sec:gainmethods}

For a stimulus $\times$ unit matrix $X$ (720 rows), let $u$ be the mean profile (unit means over
the 720 stimuli, minus their mean over units, scaled to unit norm) and $\tilde{x}_i$ row $i$
centred across units. The mean-profile alignment of stimulus $i$ is
$a_i = \langle \tilde{x}_i, u \rangle / \lVert \tilde{x}_i \rVert$; $c_{\mathrm{align}}$ is the
Pearson correlation over stimuli between model alignment (raw features) and V1 voxel alignment
($B$ = paper), averaged over subjects. Each centred pattern splits into its projection
$p_i = \langle \tilde{x}_i, u \rangle$ on the mean profile and the orthogonal remainder
$r_i = \tilde{x}_i - p_i u$. We compare ``full'' (correlation distance on $x_i$, the convention of
\citet{evalres}), ``orthogonal'' (correlation distance on $r_i$) and ``parallel''
($\lvert p_i - p_j \rvert$) RDMs; the betas are decomposed in the same way, and each model
component is compared with the same component of the brain.

For unit-centred features $X_c$ ($F$ = center) with row norms $n_i$, euclidean distances decompose
exactly as $\lVert X_{c,i} - X_{c,j} \rVert^2 = (n_i - n_j)^2 + W_{ij}$ with
$W_{ij} = 2\, n_i n_j \,(1 - \cos\theta_{ij})$. We report the norm RDM
$N_{ij} = \lvert n_i - n_j \rvert$, the unweighted angular component $A_{ij} = 1 - \cos\theta_{ij}$
and the norm-weighted angular term $W_{ij}$ separately, for features and betas alike; ``angular
component'' denotes only $A$. These names label quantities defined here; they are not claims about
physiological gain (details in Appendix~\ref{app:gainmethods}).

\subsection{Analysis plan, amendments and inference}
\label{sec:inference}

\paragraph{Pre-specified primary analysis.}
The ridge-readout analysis followed a written analysis plan with a frozen configuration
(configuration hash \texttt{206fa875cca5}), dated 20~August~2026. The plan was not publicly
registered; its dating rests on in-text date stamps and the configuration hash, and the files
were placed under version control on 3~October~2026. We therefore describe its criteria as fixed
before the analysis rather than as preregistered. The unit of inference was the training seed
($n = 5$), tested with a one-sided exact sign-flip test over all 32 sign assignments, and the
effect was to count as surviving only if it held at $g \geq 4$. A minimal detectable effect
$\Delta_{\min}$, derived from the pilot's seed-level standard deviation by simulating this test,
was compared with the reference effect reported in version~1 of \citet{evalres}
($+\PreregReferenceEffect \pm \PreregReferenceSem$ at 224\,px); the pre-specified rule labels the
outcome ``The effect disappears'' when $\Delta_{\min}$ is below the reference effect, a decision
label under the plan, not an equivalence test. Three amendments were adopted before any fit to
THINGS data and four further decisions after data had been seen, among them a post-hoc plateau
gate and a downstream-equivalence gate that replaced a failed reproduction gate for the
CPU-extracted $g = 2$ features; an implementation error in $\Delta_{\min}$ did not change the
verdict. The primary analysis was not blind (results for subject~1 at $g \in \{1, 4, 8\}$ and for
subjects~2 and~3 at $g = 1$ were known), and several pre-specified instruments were not
implemented (Appendix~\ref{app:inference}).

\paragraph{Exploratory analyses and inference.}
All analyses in Sections~\ref{sec:standardization}--\ref{sec:gaingeometry} lie outside the
pre-specified plan and were developed sequentially. For each step, predictions and decision rules
were committed to version control before the step's scripts were run, and we report every
pre-stated prediction, including the \FailedPredictions{} that were not met
(Appendix~\ref{app:failed}). The rules of later steps were written after the results of earlier
steps were known, and we chose the thresholds ourselves, so ``met'' means consistent with a
pre-stated rule, not independently confirmed; some analyses were added in response to criticism of
earlier drafts (Appendix~\ref{app:inference}). Observations without a rule are labelled
descriptive. RSA intervals are 95\% stimulus-bootstrap percentile intervals over the 720 stimuli,
for these three subjects and these networks. The 36 multiverse cells are
strongly dependent, intervals are not corrected for multiple comparisons, and counts of cells and
of ImageNet models (one set of trained weights each) are descriptive. With three subjects, the
results do not generalise to a population of observers (Table~\ref{tab:evidence}). Code,
configuration, amendments and the full commit history are available in the accompanying
repository.\footnote{\url{https://anonymous.4open.science/r/PLACEHOLDER}}
